\documentclass[10pt,a4paper]{amsart}
\usepackage[margin=19mm]{geometry}
\usepackage{amsmath,amssymb,mathtools}
\usepackage[scr=rsfs,cal=euler]{mathalfa}
\usepackage{xfrac}
\usepackage[unicode,hidelinks,bookmarks=false]{hyperref}
\newcommand{\ie}{i.e.\ }
\newcommand{\cf}{cf.\ }
\newcommand{\resp}{respectively\ }
\newcommand{\Subsection}{\subsection}
\providecommand{\nobreakdash}{\nobreak\hbox{-}}
\setlength{\emergencystretch}{2em}
\multlinegap0pt
\usepackage{bm}
\usepackage{bbm}
\usepackage{dsfont}
\usepackage{xspace}
\usepackage{cancel}
\usepackage{mathtools}
\usepackage{amsthm}
\makeatletter
\@ifundefined{theorem}{\newtheorem{theorem}{Theorem}}{}
\makeatother
\title[Further Classes of Series Involving Central Binomial Coeffic]{Further Classes of Series Involving Central Binomial Coefficients}
\author{Karl Dilcher, Christophe Vignat}
\date{Source: arXiv:2511.00109v1 (CC BY 4.0)}
\begin{document}
\maketitle

\noindent\emph{Source and licence.} Further Classes of Series Involving Central Binomial Coefficients. Version arXiv:2511.00109v1 (CC BY 4.0), distributed under the Creative Commons Attribution 4.0 International licence, as recorded in the arXiv licence metadata for this exact version and shown on the abstract page https://arxiv.org/abs/2511.00109v1. Full rights notice: licenses/arxiv-2511.00109-CC-BY-4.0.txt.\par\smallskip

\noindent\textit{Editorial scope.} Counted unit: the Theorem whose source label is \texttt{thm:5.1} in the article (source0.tex lines 684--698, source label \texttt{thm:5.1}), delivered together with the article's own complete original proof of it (source0.tex lines 700--740, one \texttt{proof} environment of 41 lines). No mathematics is written, repaired or re-derived: statement and proof are the article's own text line for line. Required context retained uncounted: 2.2 (lines 153--156); 5.1 (lines 633--636); 5.4 (lines 678--680). Every definition, display and label of those lines is retained unchanged. Omitted from this package and not redistributed: 1--152, 157--632, 637--677, 681--683, 699--699, 741--1231. Nothing inside a retained excerpt is omitted or altered: every retained body is delivered exactly as the article's lines are written, so no omission receipt of any kind accompanies this package. Citations: the retained excerpts carry no citation command at all, so no bibliography entry of the article is transcribed and no reference is redistributed. Reference adaptation: none was needed and none was made; every reference inside the delivered unit resolves inside it, so no unresolved reference or citation is produced. Printed slips kept literal: no printed word, symbol, hypothesis or number of the counted statement or of its proof was altered. Layout and class: the article's own class is \texttt{amsart} with packages including epsfig, url; the wrapper is the lane's pinned \texttt{amsart} editorial wrapper at 10pt on A4, which restates the theorem-like environments the retained text needs and adds the article's own macro definitions that the retained text needs (none), with extra wrapper packages (bm, bbm, dsfont, xspace, cancel, mathtools). The delivered numbering is the unit's own. Assets: no retained span contains a figure environment, a graphics-inclusion command, a PDF-page inclusion, a table or a TikZ picture; no image file is copied into this package, the package declares no asset dependency and no image file is redistributed with it. Licence: version arXiv:2511.00109v1 of this article is distributed under the Creative Commons Attribution 4.0 International licence, as recorded in the arXiv licence metadata for this exact version and shown on the abstract page https://arxiv.org/abs/2511.00109v1. A full rights notice is delivered at licenses/arxiv-2511.00109-CC-BY-4.0.txt. Characters: every retained excerpt is pure ASCII, so the delivered engine, XeLaTeX with its default Latin Modern text font, reports no missing character.
\begin{equation}\label{2.2}
\sum_{k=0}^\infty\frac{\binom{2k}{k}}{4^k(2k+1+z)}
=\frac{\sqrt{\pi}}{2}\cdot\frac{\Gamma(\frac{z+1}{2})}{\Gamma(\frac{z+2}{2})}.
\end{equation}

\begin{equation}\label{5.1}
\exp\left(\sum_{k=1}^\infty x_k\frac{z^k}{k!}\right)
=\sum_{n=0}^\infty B_n(x_1,\ldots,x_n)\frac{z^n}{n!};
\end{equation}

\begin{equation}\label{5.4}
\psi^{(m)}(z) = \frac{d^m}{dz^m}\psi(z)=\frac{d^{m+1}}{dz^{m+1}}\log{\Gamma(z)}.
\end{equation}

\begin{theorem}\label{thm:5.1}
For any $z\in{\mathbb C}\setminus\{-1,-3,-5,\ldots\}$ and any integer 
$p\geq 0$ we have
\begin{equation}\label{5.5}
\sum_{k=0}^\infty\frac{\binom{2k}{k}}{4^k(2k+1+z)^{p+1}}
=\frac{(-1)^p\sqrt{\pi}}{2p!}
\cdot\frac{\Gamma(\frac{z+1}{2})}{\Gamma(\frac{z+2}{2})}
\cdot B_{p}(g^{(1)}(z),\ldots,g^{(p)}(z)),
\end{equation}
where for $k\geq 1$,
\begin{equation}\label{5.6}
g^{(k)}(z)
= 2^{-k}\left(\psi^{(k-1)}(\tfrac{z+1}{2})-\psi^{(k-1)}(\tfrac{z+2}{2})\right).
\end{equation}
\end{theorem}

\begin{proof}
We fix $z\in{\mathbb C}\setminus\{-1,-3,-5,\ldots\}$ and for a complex variable
$w$ we define
\begin{equation}\label{5.7}
f(z+w):=2\sum_{k=0}^\infty\frac{\binom{2k}{k}}{4^k(2k+1+z+w)}
=\sqrt{\pi}\cdot\frac{\Gamma(\frac{z+w+1}{2})}{\Gamma(\frac{z+w+2}{2})},
\end{equation}
which is holomorphic in the domain $D:=\{w\mid w+z\not\in{\mathbb Z}_{<0}\}$.
We recall that the second equation in \eqref{5.7} comes from \eqref{2.2}.
Next we define
\begin{equation}\label{5.8}
g(z+w):=\log f(z+w)=\log\sqrt{\pi}+\log\Gamma(\frac{z+w+1}{2})
+\log\Gamma(\frac{z+w+2}{2}),
\end{equation}
and by \eqref{5.4} we have for $k\geq 1$,
\begin{equation}\label{5.9}
g^{(k)}(z)
=2^{-k}\left(\psi^{(k-1)}(\tfrac{z+1}{2})-\psi^{(k-1)}(\tfrac{z+2}{2})\right).
\end{equation}
Expanding $g(z+w)$ as a function of $w$ in a neighborhood of $0\in D$, we get 
\[
g(z+w) = g(z)+\sum_{k=1}^\infty g^{(k)}(z)\frac{w^k}{k!},
\]
and with \eqref{5.1} and \eqref{5.8} we have
\begin{align}
f(z+w)&=\exp(g(z+w))
=\sqrt{\pi}\frac{\Gamma(\frac{z+1}{2})}{\Gamma(\frac{z+2}{2})}
\cdot\exp\left(\sum_{k=1}^\infty g^{(k)}(z)\frac{w^k}{k!}\right)\label{5.10}\\
&=\sqrt{\pi}\frac{\Gamma(\frac{z+1}{2})}{\Gamma(\frac{z+2}{2})}
\sum_{n=0}^\infty B_n(g^{(1)}(z),\ldots,g^{(n)}(z))\frac{w^n}{n!}.\nonumber
\end{align}
Taking the $p$th derivative of the right-most term in \eqref{5.10} and of the
series in \eqref{5.7}, where $p$ is a non-negative integer, we obtain
\begin{equation}\label{5.11}
\sum_{k=0}^\infty\frac{\binom{2k}{k}(-1)^p p!}{4^k(2k+1+z)^{p+1}}
=\frac{\sqrt{\pi}\Gamma(\frac{z+1}{2})}{2\Gamma(\frac{z+2}{2})}
\sum_{n=p}^\infty B_n(g^{(1)}(z),\ldots,g^{(n)}(z))\frac{w^{n-p}}{(n-p)!}.
\end{equation}
Finally, with $w=0$ in \eqref{5.11} we get \eqref{5.5}, which completes the
proof.
\end{proof}

\end{document}
